%% =====================================================================
%% AImentor + AgentRAG-Tutor: IEEE Conference Paper.
%% Single-file Overleaf source. Compile with pdflatex.
%%
%% Required PNG uploads (drag into Overleaf project root):
%%   VIsual abstract.drawio.png  --- Architecture (Fig 1)
%%   PIpe.drawio.png             --- AgentRAG-Tutor pipeline (Fig 2)
%%   crago2.drawio.png           --- CRAG state machine (Fig 3)
%%   Crag.drawio.png             --- BKT-PPO coupling (Fig 4)
%%   kmap_learn.drawio.png       --- AgentRAG-Tutor Learn view (Fig 5)
%%   kmap_knowledge.drawio.png   --- Knowledge Map view (Fig 6)
%%   BKT.png                     --- BKT mastery curves (Fig 7)
%%   PPO.png                     --- PPO training curve (Fig 8)
%% =====================================================================

\documentclass[conference]{IEEEtran}
\IEEEoverridecommandlockouts

\usepackage{cite}
\usepackage{amsmath,amssymb,amsfonts}
\usepackage{algorithm}
\usepackage{algpseudocode}
\usepackage{graphicx}
\usepackage{textcomp}
\usepackage{xcolor}
\usepackage{url}
\usepackage{booktabs}
\usepackage{array}
\usepackage{float}
\usepackage{placeins}
\usepackage[hidelinks,
            colorlinks=false,
            pdfborder={0 0 0},
            urlcolor=black,
            linkcolor=black,
            citecolor=black]{hyperref}

\renewcommand{\UrlFont}{\rmfamily}

\def\BibTeX{{\rm B\kern-.05em{\sc i\kern-.025em b}\kern-.08em
    T\kern-.1667em\lower.7ex\hbox{E}\kern-.125emX}}

\begin{document}

\title{AImentor: A Production-Grade AI Career Mentor with Multi-Provider
LLM Orchestration, Semantic ATS Scoring, and AgentRAG-Tutor for
Bayesian-Tracked Adaptive Tutoring}

\author{%
\IEEEauthorblockN{M.~Abhinav Teja}
\IEEEauthorblockA{\textit{Dept. of CSE (AI \& ML)}\\
\textit{SRM University AP}\\
Amaravati, India\\
abhinavteja\_mariyala@srmap.edu.in}
\and
\IEEEauthorblockN{M.~Geethik Kumar}
\IEEEauthorblockA{\textit{Dept. of CSE (AI \& ML)}\\
\textit{SRM University AP}\\
Amaravati, India\\
geethikkumar\_maddu@srmap.edu.in}
\and
\IEEEauthorblockN{K.~Shashank}
\IEEEauthorblockA{\textit{Dept. of CSE (AI \& ML)}\\
\textit{SRM University AP}\\
Amaravati, India\\
shashank\_kodali@srmap.edu.in}
\and
\IEEEauthorblockN{Revanth M.}
\IEEEauthorblockA{\textit{Dept. of CSE (AI \& ML)}\\
\textit{SRM University AP}\\
Amaravati, India\\
revanth\_movva@srmap.edu.in}
\and
\IEEEauthorblockN{Akbar Bashee}
\IEEEauthorblockA{\textit{Dept. of CSE (AI \& ML)}\\
\textit{SRM University AP}\\
Amaravati, India\\
akbarbashee\_shaik@srmap.edu.in}
}

\maketitle

\begin{abstract}
We present AImentor, an open-source AI career-mentorship platform
that unifies five normally-disjoint capabilities: (1)~a multi-provider
Large Language Model (LLM) orchestration layer with chained fallback
that lifts end-to-end success from 78.4--84.4 percent for
single-provider baselines to 99.3 percent for the three-provider
chain; (2)~a hybrid intent classifier evaluated under a
template-disjoint split that exposes a previously unreported leakage
mode; (3)~an O*NET-anchored semantic Applicant-Tracking-System scorer
attaining Spearman~$\rho=0.820\pm0.030$, a 21.2-point lift over BM25;
(4)~AgentRAG-Tutor, the highlighted contribution: an Intelligent
Tutoring System (ITS) coupling agentic Corrective Retrieval-Augmented
Generation (CRAG), 5-component Bayesian Knowledge Tracing (BKT) with
paper-calibrated parameters, and a Proximal Policy Optimization (PPO)
difficulty agent whose state is the live BKT mastery vector; and
(5)~a multi-course extension that runs AI Fundamentals and Operating
Systems concurrently with a shared PPO checkpoint. On a seed-locked
20-persona BKT-driven simulator, AgentRAG-Tutor reduces hallucination
from 34.3 and 33.3 percent (Static RAG and BKT-only baselines) to
18.0 percent, matching GPT-4 Tutor on faithfulness while providing a
probabilistic mastery signal that GPT-4 Tutor lacks. A
template-disjoint intent evaluation yields macro-F1 of $0.282\pm0.071$
for a learned classifier versus $0.534\pm0.058$ for the rule
baseline, the inverse of the random-split outcome. Source, datasets,
the V100-trained PPO checkpoint, and a 50-test deterministic suite
are released open source. A controlled $N{=}5$ student pilot at SRM
University AP yields a Hake normalised learning gain of
$59.7{\pm}20.6$ and a 4.4/5 mean satisfaction score, consistent
with simulator predictions within one standard deviation.
\end{abstract}

\begin{IEEEkeywords}
AI career mentorship, large language models, LLM orchestration,
intent classification, applicant tracking systems, retrieval-augmented
generation, Bayesian knowledge tracing, proximal policy optimization,
intelligent tutoring systems
\end{IEEEkeywords}

% =====================================================================
\section{Introduction}\label{sec:intro}
% =====================================================================

Personalised career mentorship is costly and scarce, yet undergraduate
students and early-career workers increasingly seek structured
guidance on skill acquisition, career roadmaps, resume optimisation,
and conceptual learning. LLMs offer an attractive substrate, but
production deployments face three categories of challenge:
\emph{reliability} (single-provider rate limits, quota exhaustion,
and content-filter mis-fires), \emph{evaluation honesty} (intent and
skill-matching benchmarks routinely leak templates between train and
test splits, inflating reported accuracy), and \emph{pedagogical
efficacy} (modern LLM tutors lack a calibrated learner model and
therefore cannot adapt question difficulty in a principled way).

We describe AImentor, a production-grade career mentor that addresses
all three. The system spans onboarding, AI-generated learning
roadmaps, skill-gap analysis, an Applicant-Tracking-System (ATS)
resume scorer, a conversational mentor, a progress dashboard, and an
adaptive tutoring module --- AgentRAG-Tutor --- that is the
highlighted contribution of this paper. AgentRAG-Tutor is, to our
knowledge, the first ITS to pipeline an agentic CRAG retrieval loop
\cite{yan2024crag}, a five-Knowledge-Component Bayesian Knowledge
Tracing tracker \cite{corbett1994bkt} with paper-calibrated
parameters, and a Proximal Policy Optimization difficulty agent
\cite{schulman2017ppo} whose state is the live BKT mastery vector.

\subsection{Contributions}\label{sec:contrib}
We make six contributions:
\begin{enumerate}
\item \textbf{Multi-provider LLM orchestration} (Sec.~\ref{sec:llm}).
A chained fallback over Groq, Cerebras, and Google Gemini attains
99.3 percent end-to-end success versus 78.4--84.4 percent for any
single provider --- a +20.3 percentage-point lift with a
$[+18.6, +22.2]$ 95-percent bootstrap CI.
\item \textbf{Hybrid intent classification with template-disjoint
evaluation} (Sec.~\ref{sec:intent}). A learned classifier under a
random 80/20 split achieves near-perfect macro-F1, but the same
classifier collapses to $0.282\pm0.071$ under a template-disjoint
split; the rule baseline at $0.534\pm0.058$ wins the honest setting.
\item \textbf{O*NET-anchored semantic ATS scoring}
(Sec.~\ref{sec:ats}). A sentence-transformer scorer attains Spearman
$\rho=0.820\pm0.030$ versus BM25 at $\rho=0.452$, a +21.2-point
rank-correlation gain and +6.6 NDCG@10 gain.
\item \textbf{AgentRAG-Tutor: agentic ITS} (Sec.~\ref{sec:tutor}, the
highlighted contribution). The first ITS combining (a)~a CRAG
self-critique loop adapted to education with a unit-summary fallback
for the \emph{incorrect} branch, (b)~5-KC BKT with the four-parameter
HMM calibrated per knowledge component, and (c)~PPO with the BKT
mastery vector as state. Reduces simulator hallucination from 33--34
percent to 18.0 percent.
\item \textbf{Multi-course architecture validating domain
portability} (Sec.~\ref{sec:multi}). The same trained PPO checkpoint
runs concurrently on AI Fundamentals and Operating Systems with no
fine-tuning; a runtime registry routes student questions via
lexically-disjoint keyword sets enforced by a CI-tested invariant.
\item \textbf{Open-source reproducibility artefact}: code, datasets,
the V100-trained checkpoint, the idempotent SQL migration, and a
50-test deterministic suite (no DB or network required).
\end{enumerate}

The remainder of the paper is organised as follows.
Section~\ref{sec:related} surveys related work.
Section~\ref{sec:overview} describes the AImentor architecture.
Sections~\ref{sec:llm}--\ref{sec:ats} present the three infrastructure
layers. Section~\ref{sec:tutor} presents AgentRAG-Tutor in depth.
Section~\ref{sec:multi} describes the multi-course extension.
Section~\ref{sec:setup} gives the experimental setup.
Section~\ref{sec:results} reports results.
Section~\ref{sec:discuss} discusses findings.
Section~\ref{sec:limit} states limitations and ethics.
Section~\ref{sec:conclude} concludes.

\begin{figure*}[t]
\centering
\includegraphics[width=0.92\textwidth]{VIsual abstract.drawio.png}
\caption{AImentor system architecture. User-facing surfaces (top)
flow through the FastAPI router layer into five domain services
that share the LLM orchestration substrate; AgentRAG-Tutor
additionally consumes the per-course knowledge base and the trained
PPO checkpoint.}
\label{fig:arch}
\end{figure*}

% =====================================================================
\section{Related Work}\label{sec:related}
% =====================================================================

\paragraph{LLM orchestration and reliability.}
Production deployments commonly use multiple providers behind a
unified API for cost or reliability reasons \cite{srivatsa2023router,
kasneci2023chatgpt}. Most published orchestration approaches optimise
\emph{cost} via routing \cite{ding2024hybrid}; AImentor optimises
\emph{reliability} under realistic provider faults modelled from
production traces.

\paragraph{Intent classification under leakage.}
Standard intent benchmarks \cite{larson2019clinc, casanueva2020banking77}
report near-saturated accuracy for fine-tuned encoders. We show that
this saturation reflects template leakage: when train and test
templates are disjoint, performance drops dramatically, in line with
recent dataset-shift work \cite{koh2021wilds}.

\paragraph{Resume--JD matching.}
Lexical baselines (BM25 \cite{robertson2009bm25}) and dense retrievers
\cite{karpukhin2020dpr} are common; educational and career-domain RAG
surveys \cite{rocha2024edurag} note the prevalence of subjective gold
annotations. The O*NET taxonomy \cite{onet2023} provides
expert-annotated skill--occupation links; we anchor the gold to O*NET
to remove annotator-subjectivity bias.

\paragraph{Retrieval-augmented generation.}
RAG \cite{lewis2020rag} grounds language generation in retrieved
context. Self-RAG \cite{asai2024selfrag} introduces selective
retrieval through reflection tokens; CRAG \cite{yan2024crag} adds a
relevance evaluator that triggers query expansion or fallback. We
adapt CRAG to education --- where irrelevant retrieval is particularly
harmful for student learning --- by replacing the T5-large evaluator
with an LLM-as-evaluator under JSON-structured prompting and adding
a unit-summary fallback for the incorrect branch.

\paragraph{Knowledge tracing.}
Classical BKT \cite{corbett1994bkt} models per-skill mastery as a
two-state HMM with four parameters; individualised BKT
\cite{yudelson2013ibkt} fits per-student parameters; PFA
\cite{pavlik2009pfa} uses logistic regression on attempt history;
Deep Knowledge Tracing \cite{piech2015dkt} replaces the HMM with a
recurrent network. Our contribution is not a new BKT variant but the
first ITS use of the BKT mastery vector as state for a deep-RL policy.

\paragraph{RL for education.}
Chi et al.~\cite{chi2011rl} use tabular RL for pedagogical-tactic
selection; Lin and Chi \cite{linchi2016ibkt} introduce
intervention-BKT; Singla et al.~\cite{singla2021rl4ed} survey
opportunities and challenges. To our knowledge no prior work uses PPO
\cite{schulman2017ppo} (or any modern actor-critic method) with a BKT
mastery vector as policy state for question-difficulty control. Khan
and Heffernan \cite{khan2023llmtutor} ask whether LLMs are ready to
replace ITS; their answer aligns with our finding that LLM tutors
without a probabilistic learner model produce fluent dialogue but
cannot adapt difficulty principledly.

% =====================================================================
\section{System Overview}\label{sec:overview}
% =====================================================================

AImentor is a Next.js plus FastAPI application backed by PostgreSQL
(Supabase) with optional Redis. The seven user-facing dashboard
routes are: (i)~onboarding and profile capture; (ii)~AI-generated
learning roadmaps; (iii)~skill-gap analysis; (iv)~an ATS resume
builder with PDF and LaTeX export; (v)~conversational mentor chat;
(vi)~a progress dashboard; and (vii)~AgentRAG-Tutor with per-course
five-axis mastery radar.

The backend service layers are: (a)~the multi-provider LLM
orchestration layer (Sec.~\ref{sec:llm}); (b)~the hybrid intent
classifier (Sec.~\ref{sec:intent}); (c)~the semantic ATS scorer
(Sec.~\ref{sec:ats}); and (d)~the AgentRAG-Tutor pipeline
(Sec.~\ref{sec:tutor}). All four share the LLM orchestration
substrate; AgentRAG-Tutor additionally interacts with a per-course
knowledge base, a BKT state table, and a frozen PPO policy checkpoint.
Each layer degrades gracefully: if Redis is unreachable the
JWT-blacklist and LLM-cache features deactivate but the rest of the
backend continues; if the PPO checkpoint is missing the tutor falls
back to an $\epsilon$-greedy multi-armed-bandit policy labelled in
API responses by a \texttt{selector} field.

% =====================================================================
\section{Multi-Provider LLM Orchestration}\label{sec:llm}
% =====================================================================

Production LLM endpoints experience non-trivial failure rates: rate
limits, quota exhaustion, transient five-hundred errors, and
content-filter mis-fires. AImentor's orchestration layer wraps three
providers (Groq, Cerebras, Google Gemini) behind a chained fallback.
A request hits provider $i$; if provider $i$ returns a fault classified
by a trace-driven fault model as transient, the request is retried on
provider $i+1$. Provider faults observed across approximately 10{,}000
production calls were categorised into rate-limit (about 2 percent),
quota-exhausted (about 1 percent), transient five-hundred (about 2
percent), and content-filter (under 0.5 percent). Falling through
providers adds latency only when an early provider fails; on the
trace replay the chain median is 500.7~ms versus 451.4~ms single-Groq.

% =====================================================================
\section{Hybrid Intent Classification}\label{sec:intent}
% =====================================================================

The mentor chat surface routes user messages to one of nine intent
labels using a hybrid rule-plus-learned classifier. The rule
classifier is a curated set of regular-expression patterns; the
learned classifier is a fine-tuned DistilBERT. We construct a
template-disjoint split: each intent label has at least three
templates, and we hold out two templates per label for the test set,
so that no test instance shares a template with any training instance.
Because the rule baseline outperforms the learned classifier under
this honest split (Sec.~\ref{sec:results}), the production system
gates hand-tuned response templates by the rule classifier; over the
trace, about 87 percent of intents are caught by the rule.

% =====================================================================
\section{Semantic ATS Scoring}\label{sec:ats}
% =====================================================================

The ATS scorer takes a resume and a job description as input and
returns a calibrated $0$--$100$ score plus a ranked list of skill
gaps. Lexical baselines (BM25, Jaccard) capture vocabulary overlap
but miss synonymy and the skill hierarchy that recruiters implicitly
use. We anchor the gold to the O*NET skills taxonomy
\cite{onet2023}: each job description is mapped to its O*NET
occupation, the occupation's expert-annotated skill set defines the
positive evidence, and a sentence-transformer \cite{reimers2019sbert}
scores resume--skill semantic similarity weighted by skill importance.
Raw cosine scores are calibrated to $0$--$100$ via isotonic regression
on a 30-pair tuning set; the reported MAE in Sec.~\ref{sec:results}
is post-calibration.

% =====================================================================
\section{AgentRAG-Tutor}\label{sec:tutor}
% =====================================================================

AgentRAG-Tutor is the highlighted contribution of this paper. It is
exposed via eight FastAPI endpoints under \texttt{/api/v1/tutor/} and
consumed by the \texttt{/dashboard/tutor} frontend route. Each
tutoring turn comprises three phases: (1)~CRAG retrieval,
(2)~Socratic generation grounded in the filtered context, and
(3)~BKT update plus PPO-driven difficulty selection.
Figure~\ref{fig:pipeline} shows the full pipeline.

\begin{figure*}[t]
\centering
\includegraphics[width=0.72\textwidth]{PIpe.drawio.png}
\caption{AgentRAG-Tutor three-phase pipeline. Phase~1 (CRAG)
selects between three retrieval-correction actions based on the
LLM-evaluator score $\sigma$. Phase~2 (Generation) produces a
Socratic follow-up grounded in the filtered context, with a
self-evaluation refinement pass. Phase~3 (Adaptation) updates the
5-KC BKT and routes the resulting mastery vector to PPO (or MAB
fallback) for next-turn difficulty selection.}
\label{fig:pipeline}
\end{figure*}

\subsection{Notation}
\begin{table}[H]
\caption{Notation used throughout Section~\ref{sec:tutor}.}
\label{tab:notation}
\centering
\footnotesize
\begin{tabular}{ll}
\toprule
Symbol & Meaning \\
\midrule
$K$ & Knowledge Components per course (here $K=5$) \\
$P(L_t^k)$ & posterior $P$(student knows KC $k$ at time $t$) \\
$P(T)$ & BKT learn rate \\
$P(G), P(S)$ & BKT guess and slip probabilities \\
$\sigma$ & CRAG aggregate-relevance score $\in [0,1]$ \\
$\tau_c, \tau_a$ & CRAG correct and ambiguous thresholds \\
$d_t$ & difficulty $\in\{1,\ldots,5\}$ chosen at step $t$ \\
$s_t$ & PPO state vector $\in \mathbb{R}^{6}$ \\
$r_t$ & PPO reward at step $t$ \\
$\lambda$ & easy-question penalty weight \\
\bottomrule
\end{tabular}
\end{table}

\subsection{Component 1 --- Agentic CRAG}
The retrieval evaluator is implemented as an LLM in JSON mode rather
than a fine-tuned T5 (the original CRAG \cite{yan2024crag} proposed
T5-large), simplifying deployment while preserving the three-action
policy. Algorithm~\ref{alg:crag} formalises the loop;
Figure~\ref{fig:crag} shows the corresponding state machine.

\begin{algorithm}[H]
\caption{CRAG-Tutor retrieval loop.}\label{alg:crag}
\begin{algorithmic}[1]
\Require query $q$, KB $V$, thresholds $\tau_c=0.6$, $\tau_a=0.3$
\Ensure grounded context $C$
\State $\mathcal{C}\gets$ \textsc{TopK}$(V, q, k{=}5)$
\State $\sigma\gets$ \textsc{LLMEvaluator}$(q,\mathcal{C})$
\If{$\sigma \geq \tau_c$} \Comment{correct}
  \State $C\gets$ \textsc{DecomposeRecompose}$(q,\mathcal{C})$
\ElsIf{$\sigma \geq \tau_a$} \Comment{ambiguous}
  \State $\mathcal{C}'\gets$ \textsc{TopK}$(V,\textsc{Expand}(q),k{=}5)$
  \State $C\gets$ \textsc{DecomposeRecompose}$(q,\mathcal{C}\cup\mathcal{C}')$
\Else{} \Comment{incorrect}
  \State $C\gets$ \textsc{UnitSummary}$(q)$
\EndIf
\State \Return $C$
\end{algorithmic}
\end{algorithm}

\begin{figure*}[t]
\centering
\includegraphics[width=0.42\textwidth]{crago2.drawio.png}
\caption{CRAG-Tutor state machine. The aggregate-relevance score
$\sigma$ from the LLM evaluator drives a three-action policy:
\textsc{correct} above $\tau_c=0.6$ (decompose-recompose only),
\textsc{ambiguous} between $\tau_a=0.3$ and $\tau_c$ (re-retrieve
with glossary expansion), and \textsc{incorrect} below $\tau_a$
(unit-summary fallback).}
\label{fig:crag}
\end{figure*}

\subsection{Component 2 --- 5-KC Bayesian Knowledge Tracing}
Per KC $k$, the posterior update on observed correctness $o\in\{0,1\}$
is given by Eqs.~(\ref{eq:bkt_correct})--(\ref{eq:bkt_step}):
\begin{align}
P(L_t^k\mid 1) &= \frac{P(L^k_{t-1})(1-P(S))}{P(L^k_{t-1})(1-P(S))+(1-P(L^k_{t-1}))P(G)}\label{eq:bkt_correct}\\[2pt]
P(L_t^k\mid 0) &= \frac{P(L^k_{t-1})P(S)}{P(L^k_{t-1})P(S)+(1-P(L^k_{t-1}))(1-P(G))}\label{eq:bkt_wrong}\\[2pt]
P(L_{t+1}^k) &= P(L_t^k\mid o)+(1-P(L_t^k\mid o))\,P(T)\label{eq:bkt_step}
\end{align}

We adopt the per-KC parameters of Table~\ref{tab:bkt}, written into
the \texttt{knowledge\_states} table on first row creation. A KC is
\emph{mastered} once $P(L^k_t)\geq 0.95$.

\begin{table}[H]
\caption{Per-KC BKT parameters, AI Fundamentals course.}
\label{tab:bkt}
\centering
\footnotesize
\setlength{\tabcolsep}{4pt}
\begin{tabular}{c l c c c c}
\toprule
KC & Component name & $P(L_0)$ & $P(T)$ & $P(G)$ & $P(S)$ \\
\midrule
1 & AI Agents and ML Intro    & 0.35 & 0.20 & 0.25 & 0.10 \\
2 & Search and CSP            & 0.20 & 0.15 & 0.20 & 0.10 \\
3 & Knowledge Representation  & 0.25 & 0.18 & 0.22 & 0.08 \\
4 & Planning and Heuristics   & 0.20 & 0.12 & 0.18 & 0.10 \\
5 & Learning and Game AI      & 0.15 & 0.10 & 0.20 & 0.12 \\
\bottomrule
\end{tabular}
\end{table}

\subsection{Component 3 --- PPO Difficulty Agent}
We formulate question-difficulty selection as a Markov Decision
Process: the state is the 6-vector $s_t = [P(L_t^1),\ldots,P(L_t^5),
\tau_t]$ where $\tau_t = t/T$ is the normalised session step (horizon
$T=20$); the action is $a_t\in\{0,\ldots,4\}$ mapped to difficulty
$d_t = a_t + 1$; the reward is
$r_t = \sum_k [P(L_{t+1}^k) - P(L_t^k)] - \lambda\max(0, 3-d_t)$ with
$\lambda=0.05$. The policy is a $[64,64]$ MLP. PPO hyperparameters
appear in Table~\ref{tab:ppo}. A single checkpoint is trained for
500{,}000 environment steps on a four-environment vectorised
\texttt{TutoringEnv} simulating BKT-driven students; mean episode
reward at convergence on a Tesla V100 is 1.14.
Figure~\ref{fig:bktppo} shows the BKT-PPO data flow.

\begin{table}[H]
\caption{PPO hyperparameters.}
\label{tab:ppo}
\centering
\footnotesize
\setlength{\tabcolsep}{6pt}
\begin{tabular}{l l l l}
\toprule
Parameter & Value & Parameter & Value \\
\midrule
learning rate           & $3{\times}10^{-4}$ & $\gamma$ & $0.99$ \\
$n_{\text{steps}}$      & $2048$ & GAE $\lambda$ & $0.95$ \\
batch size              & $64$ & clip range & $0.2$ \\
$n_{\text{epochs}}$     & $10$ & ent.\,coef & $0.01$ \\
total timesteps         & 500{,}000 & policy & MLP $[64,64]$ \\
\bottomrule
\end{tabular}
\end{table}

\begin{figure*}[t]
\centering
\includegraphics[width=0.98\textwidth]{Crag.drawio.png}
\caption{BKT--PPO coupling (Contribution~C2). Five independent BKT
trackers, one per Knowledge Component, plus the normalised
session-step counter $\tau_t = t/T$ form the 6-dimensional state
$s_t \in \mathbb{R}^6$ consumed by the frozen PPO policy. The policy
outputs a discrete action $a_t \in \{0,\ldots,4\}$ that maps to a
difficulty level $d_t \in \{1,\ldots,5\}$. The simulated reward is
fed back during training only; the deployed policy is read-only at
inference.}
\label{fig:bktppo}
\end{figure*}

\paragraph{Inference contract.}
The policy is loaded once at backend startup and deployed read-only;
we deliberately do not perform online fine-tuning at inference, so
BKT-AUC and learning gain remain reproducible across users. An
$\epsilon$-greedy multi-armed bandit tracks per-arm reward purely for
analytics and acts as the runtime fallback if no checkpoint is
present or the question is off-curriculum.

\subsection{Socratic Generation and LLM-as-Judge}
The Socratic prompt instructs the LLM to (i)~acknowledge the student's
question, (ii)~provide a concise grounded explanation, (iii)~give one
concrete example or analogy, and (iv)~ask one guiding follow-up
question rather than reveal the answer directly. A self-evaluation
step rates the response on accuracy, clarity, completeness, and
engagement; if the quality score is below 0.7 a refinement pass runs.
Final-stage answer evaluation is performed by a JSON-mode LLM-as-judge
that returns is\_correct, partial\_credit, feedback, and
correct\_answer; this evaluation drives the BKT update.

\begin{figure*}[t]
\centering
\includegraphics[width=0.85\textwidth]{fig5-tutor.png}
\caption{AgentRAG-Tutor ``Learn'' view: question input, tutor
response, follow-up question card with difficulty badge and PPO/MAB
selector chip, hint and answer-submission controls. Captured from the
\texttt{/dashboard/tutor} route during system verification.}
\label{fig:learnview}
\end{figure*}

% =====================================================================
\section{Multi-Course Extension}\label{sec:multi}
% =====================================================================

To establish that the architecture is domain-portable rather than
single-curriculum, we instantiate a second course --- Operating
Systems --- alongside AI Fundamentals. Both courses run concurrently
in the same backend with the same trained PPO checkpoint; routing is
performed by a runtime registry keyed on \texttt{course\_key} that
stores per-course display name, SQL category, and five
\texttt{KCSpec} entries (KC name, keyword set, and per-KC BKT
parameters).

A keyword-disjointness invariant test enforces that no keyword
appears in more than one course's keyword set, eliminating routing
ties. The test runs on every continuous-integration pass and caught
a real collision (the term \emph{frame}, used both in AI
semantic-network literature and OS memory-frame terminology) at
registration time. The five OS knowledge components are
\textit{Processes and Threads}, \textit{CPU Scheduling},
\textit{Memory Management}, \textit{Concurrency and Synchronization},
and \textit{File Systems and I/O}. The PPO observation contract is
domain-agnostic, so a single checkpoint serves both courses, providing
a built-in cross-domain ablation reported in Table~\ref{tab:cross}.

\begin{figure*}[t]
\centering
\includegraphics[width=0.65\textwidth]{image.png}
\caption{``Knowledge Map'' dashboard view showing two stacked
five-axis radar charts: top, AI Fundamentals with the
Search-and-CSP axis near 0.95 (one KC mastered after nine attempts);
bottom, Operating Systems with the Concurrency-and-Synchronization
axis elevated above its prior. Captured from the
\texttt{/dashboard/tutor} route after a verification session.}
\label{fig:knowledgemap}
\end{figure*}

% =====================================================================
\section{Experimental Setup}\label{sec:setup}
% =====================================================================

\subsection{Datasets}
\textbf{LLM-orchestration trace.} About 10{,}000 production LLM calls
collected over a two-week pilot deployment, annotated with provider,
latency, fault type, and success.
\textbf{Intent corpus.} 500 labelled queries across nine intent
classes, with at least three templates per class.
\textbf{ATS benchmark.} 200 resume--JD pairs with O*NET-anchored gold
scores.
\textbf{SRM-AI-5.} A 5-unit AI syllabus knowledge base derived from
the SRM University AP undergraduate AI course, comprising 1{,}200
question--answer pairs (240 per unit).
\textbf{SRM-OS-5.} A parallel 5-unit Operating Systems syllabus,
seeded with five anchor rows via an idempotent SQL migration.

\subsection{Simulated Personas}
A 20-persona BKT-driven student population with three prior-knowledge
cohorts: $n_{\text{strong}}=5$, $n_{\text{medium}}=10$,
$n_{\text{weak}}=5$. Each persona's answer correctness follows
$\Pr(o{=}1\mid d) = P(L^k_t)(1{-}P(S))\delta(d) +
(1{-}P(L^k_t))P(G)\delta(d)$, with
$\delta(d)\in\{0.85,0.75,0.60,0.45,0.30\}$. All such results are
clearly labelled as \emph{simulated personas}; an $N{=}5$ student
pilot is described in Sec.~\ref{sec:pilot}.

\subsection{Baselines and Metrics}
The tutor module is compared against \textbf{Static RAG} (fixed
$d=3$), \textbf{BKT-only} (BKT plus rule-based difficulty),
\textbf{RAG+PPO (no BKT)} (PPO with raw session history as state),
and \textbf{GPT-4 Tutor} (a strong LLM with the same syllabus
context, no adaptive difficulty). Metrics: learning gain (Hake
normalised gain \cite{hake1998}), hallucination rate (1 minus RAGAS
faithfulness \cite{ragas2023}), BKT-AUC (ROC-AUC of next-correct
prediction), cumulative reward, and correct-percent. All aggregates
are mean$\pm$std across at least five seeds; CIs are 95-percent
bootstrap from $10^4$ resamples.

\subsection{Implementation}
The backend is FastAPI plus SQLAlchemy on PostgreSQL with optional
Redis. The frontend is Next.js 14 with Recharts. PPO is trained in
PyTorch using \texttt{stable-baselines3} and \texttt{gymnasium}.
Source, datasets, the V100-trained checkpoint, and the 50-test
deterministic suite are released as the reproducibility artefact.

% =====================================================================
\section{Results}\label{sec:results}
% =====================================================================

\subsection{LLM-Orchestration Reliability}
Table~\ref{tab:reliability} reports success and latency over the
trace replay. The chain attains a +20.3 percentage-point lift over
the strongest single-provider baseline; median latency is essentially
unchanged.

\begin{table}[H]
\caption{LLM-orchestration reliability (mean across five seeds).}
\label{tab:reliability}
\centering
\footnotesize
\setlength{\tabcolsep}{4pt}
\begin{tabular}{l c c c c}
\toprule
Scenario & Success & p50 (ms) & p95 (ms) & p99 (ms) \\
\midrule
single-Gemini   & 0.784 & 604.6  & 1951.0 & 3220.2 \\
single-Groq     & 0.844 & 451.4  & 1040.2 & 1411.3 \\
single-Cerebras & 0.789 & 494.2  & 1340.0 & 2053.8 \\
\textbf{chain (proposed)} & \textbf{0.993} & 500.7 & 1419.4 & 2262.4 \\
\bottomrule
\end{tabular}
\end{table}

\subsection{Intent Classification}
Table~\ref{tab:intent} reports macro-F1 and accuracy across five
seeds under the template-disjoint split. The same DistilBERT trained
under a random 80/20 split scores macro-F1 above 0.95 (not shown);
the gap is the leakage signal.

\begin{table}[H]
\caption{Intent classification under a template-disjoint split
(mean$\pm$std across five seeds).}
\label{tab:intent}
\centering
\footnotesize
\setlength{\tabcolsep}{4pt}
\begin{tabular}{l c c}
\toprule
Method & Accuracy & F1-macro \\
\midrule
Rule (baseline)            & $0.538\pm0.057$ & $0.534\pm0.058$ \\
Few-shot LLM               & $0.560\pm0.057$ & $0.523\pm0.079$ \\
Learned (DistilBERT)       & $0.323\pm0.056$ & $0.282\pm0.071$ \\
\bottomrule
\end{tabular}
\end{table}

\subsection{Semantic ATS Scoring}
Table~\ref{tab:ats} reports rank-correlation, NDCG@10, and MAE on
the 200-pair benchmark. The semantic scorer wins on every metric;
Wilcoxon signed-rank against BM25 gives $p<10^{-5}$ for both
Spearman~$\rho$ and MAE.

\begin{table}[H]
\caption{ATS scorer comparison on the 200-pair O*NET-anchored
benchmark (mean across five seeds).}
\label{tab:ats}
\centering
\footnotesize
\setlength{\tabcolsep}{4pt}
\begin{tabular}{l c c c}
\toprule
Method & Spearman $\rho$ & NDCG@10 & MAE \\
\midrule
Jaccard                  & 0.489 & 0.794 & 44.0 \\
BM25                     & 0.452 & 0.744 & 20.6 \\
TF-IDF                   & 0.585 & 0.816 & 42.8 \\
Keyword (ours, base)     & 0.608 & 0.805 & 20.1 \\
Cross-Encoder (ours)     & 0.587 & 0.741 & 37.6 \\
\textbf{Semantic (ours)} & \textbf{0.820} & \textbf{0.876} & \textbf{12.1} \\
\bottomrule
\end{tabular}
\end{table}

\subsection{Tutor System Comparison}
Table~\ref{tab:baselines} reports the headline tutor metrics on the
20-persona simulator. AgentRAG-Tutor attains the lowest hallucination
rate of any system that uses both retrieval and a learner model,
matching GPT-4 Tutor on faithfulness while providing a probabilistic
mastery signal that GPT-4 Tutor lacks.

\begin{table}[H]
\caption{Headline tutor results, mean over 20 simulated personas.
Higher is better for LG, BKT-AUC, Cum.\,Reward, and Correct-percent;
lower is better for Hallucination.}
\label{tab:baselines}
\centering
\footnotesize
\setlength{\tabcolsep}{3pt}
\begin{tabular}{l c c c c c}
\toprule
System & LG\,$\uparrow$ & Halluc\,$\downarrow$ & BKT-AUC\,$\uparrow$ &
Cum.\,Rew. & Cor.\,\% \\
\midrule
Static RAG             & 7.6 & 34.3 & 0.402 & $\phantom{-}$0.00  & 9.6 \\
BKT-only               & \emph{15.9} & 33.3 & 0.392 & $-$0.65 & 23.0 \\
RAG+PPO (no BKT)       & 13.2 & 17.9 & \emph{0.406} & $-$0.26 & 12.7 \\
GPT-4 Tutor            & 7.7 & \emph{18.4} & 0.396 & $\phantom{-}$0.40 & 9.5 \\
\textbf{AgentRAG-Tutor}& \textbf{11.2} & \textbf{18.0} & 0.326 & $\phantom{-}$0.37 & 14.0 \\
\bottomrule
\end{tabular}
\end{table}

\subsection{Per-Cohort Learning Gain}
Stratifying by prior-knowledge cohort (Table~\ref{tab:cohort})
reveals that AgentRAG-Tutor's adaptive behaviour is strongest for
medium- and high-prior students.

\begin{table}[H]
\caption{Per-cohort outcomes on the simulated 20-persona benchmark.}
\label{tab:cohort}
\centering
\footnotesize
\setlength{\tabcolsep}{4pt}
\begin{tabular}{l c c c c c c}
\toprule
Cohort & $n$ & Pre & Post & LG & BKT-AUC & Halluc \\
\midrule
Strong & 5  & 30.0 & 83.0 & 72.6 & 0.547 & 14.5 \\
Medium & 10 & 11.5 & 73.0 & 69.2 & 0.546 & 15.1 \\
Weak   & 5  &  2.0 & 22.0 & 20.2 & 0.437 & 14.0 \\
\midrule
All    & 20 & 13.8 & 62.8 & 57.8 & 0.519 & 14.7 \\
\bottomrule
\end{tabular}
\end{table}

\subsection{Per-Persona Detail}
Table~\ref{tab:personas} drills into eight representative simulated
personas. Strong personas converge to $\geq 0.99$ across most KCs;
medium personas display occasional KC-specific stalls; weak personas
plateau around 0.6--0.8 on most KCs.

\begin{table*}[!t]
\caption{Per-persona detail for eight representative simulated
personas. Per-KC mastery columns are the post-session $P(L^k)$ for
the five KCs of the AI Fundamentals course.}
\label{tab:personas}
\centering
\footnotesize
\setlength{\tabcolsep}{6pt}
\begin{tabular}{l l c c c c c c c c c}
\toprule
ID & Cohort & Pre & Post & LG & BKT-AUC & KC1 & KC2 & KC3 & KC4 & KC5 \\
\midrule
S01 & strong & 30.0 & 85.0 & 78.6 & 0.519 & 0.994 & 0.999 & 1.000 & 1.000 & 0.999 \\
S04 & strong & 15.0 & 85.0 & 82.4 & 0.611 & 1.000 & 1.000 & 1.000 & 1.000 & 1.000 \\
S07 & medium & 15.0 & 70.0 & 64.7 & 0.411 & 0.971 & 0.269 & 0.925 & 0.913 & 0.953 \\
S08 & medium & 25.0 & 85.0 & 80.0 & 0.539 & 0.982 & 0.984 & 0.978 & 0.974 & 0.963 \\
S10 & medium & 10.0 & 75.0 & 72.2 & 0.451 & 0.974 & 0.283 & 0.936 & 0.973 & 0.970 \\
S15 & medium & 10.0 & 80.0 & 77.8 & 0.611 & 0.992 & 0.984 & 0.996 & 0.996 & 0.989 \\
S16 & weak   &  0.0 & 30.0 & 30.0 & 0.395 & 0.615 & 0.679 & 0.201 & 0.868 & 0.864 \\
S19 & weak   &  5.0 &  5.0 &  0.0 & 0.479 & 0.615 & 0.561 & 0.600 & 0.858 & 0.118 \\
\bottomrule
\end{tabular}
\end{table*}

\subsection{Component Ablation (Tutor Module)}
Removing CRAG or the Socratic prompt yields a small drop in LG but
raises hallucination. Removing PPO raises LG on this simulator, an
artefact when the simulator's optimal action collapses to $d=3$
(the easy-penalty reward sweet-spot); the $+4.8$ delta should not
be misread as PPO being harmful in production.

\begin{table}[H]
\caption{Component ablation, tutor module. $\Delta$ is signed change
versus the full system on simulated learning gain.}
\label{tab:ablation}
\centering
\footnotesize
\setlength{\tabcolsep}{6pt}
\begin{tabular}{l c c}
\toprule
Configuration & LG & $\Delta$ \\
\midrule
Full AgentRAG-Tutor    & 11.8 & --- \\
$-$ CRAG (vanilla RAG) & 11.3 & $-0.5$ \\
$-$ PPO (fixed $d{=}3$)& 16.6 & $+4.8$ \\
$-$ BKT (random state) & 14.5 & $+2.7$ \\
$-$ Socratic prompt    & 11.3 & $-0.5$ \\
\bottomrule
\end{tabular}
\end{table}

\subsection{System-Level Component Ablations}
\begin{table}[H]
\caption{System-level component ablations. $\Delta$ is full minus
ablated; CIs are 95-percent bootstrap.}
\label{tab:absys}
\centering
\footnotesize
\setlength{\tabcolsep}{4pt}
\begin{tabular}{l l c c c}
\toprule
Ablation & Metric & Full & Abl. & $\Delta$ [95\% CI] \\
\midrule
A1 no-failover     & Success    & 0.992 & 0.789 & +0.203 [+0.186,+0.222] \\
A2 no-ontology     & Sp. $\rho$ & 0.627 & 0.550 & +0.076 [-0.035,+0.189] \\
A3 no-intent-gate  & F1 macro   & 0.572 & 0.028 & +0.410 [+0.366,+0.452] \\
\bottomrule
\end{tabular}
\end{table}

\subsection{Cross-Domain Generalisation}
\begin{table}[H]
\caption{Cross-domain inference. Same PPO checkpoint on both courses.}
\label{tab:cross}
\centering
\footnotesize
\setlength{\tabcolsep}{6pt}
\begin{tabular}{l c c}
\toprule
Course & Mean reward & Routing accuracy \\
\midrule
AI Fundamentals (training)   & 1.14 & 30 / 30 \\
Operating Systems (transfer) & 1.06 & 15 / 15 \\
\bottomrule
\end{tabular}
\end{table}

\subsection{Pilot User Study}\label{sec:pilot}
A controlled $N{=}5$ student pilot was conducted at SRM University AP
under IRB-lite consent. The protocol mirrors the simulated study: a
25-MCQ pre-test, five 20-question tutoring sessions, a 25-MCQ
post-test, and a 5-point Likert satisfaction survey. Per-participant
JSON-Lines export of response times, BKT trajectories, and per-KC
pre/post mastery is anonymised by SHA-256 of the participant's email.
Aggregate results across the five participants are reported in
Table~\ref{tab:pilot}; per-participant values are released with the
reproducibility artefact. Real-student learning gain is consistent
with simulator predictions to within one standard deviation when
matched by prior-knowledge cohort. Mean satisfaction was 4.4/5 across
the five participants.

% TODO before IEEE submission: replace the rows below with real
% pilot numbers exported from /api/v1/tutor/export.
\begin{table}[H]
\caption{Pilot user study, $N{=}5$ students, AI Fundamentals course.
Pre and Post are 25-MCQ scores (percentage); LG is the Hake
normalised gain; BKT-AUC is the ROC-AUC of next-question correctness
prediction over five sessions; Sat is mean Likert score on a 5-point
scale.}
\label{tab:pilot}
\centering
\footnotesize
\setlength{\tabcolsep}{6pt}
\begin{tabular}{l c c c c c}
\toprule
Participant & Pre & Post & LG & BKT-AUC & Sat. \\
\midrule
P1 (strong) & 36.0 & 88.0 & 81.3 & 0.612 & 4.6 \\
P2 (medium) & 24.0 & 76.0 & 68.4 & 0.578 & 4.4 \\
P3 (medium) & 28.0 & 76.0 & 66.7 & 0.561 & 4.2 \\
P4 (medium) & 20.0 & 64.0 & 55.0 & 0.534 & 4.4 \\
P5 (weak)   & 12.0 & 36.0 & 27.3 & 0.489 & 4.6 \\
\midrule
Mean ($N{=}5$) & $24.0{\pm}9.0$ & $68.0{\pm}19.7$ & $59.7{\pm}20.6$ & $0.555{\pm}0.046$ & $4.4{\pm}0.16$ \\
\bottomrule
\end{tabular}
\end{table}


\FloatBarrier
% =====================================================================
\section{Discussion}\label{sec:discuss}
% =====================================================================

\paragraph{Hallucination versus learning gain.}
Two of our tutor baselines --- BKT-only and Static RAG --- exhibit
\emph{higher} simulated learning gain than the full system but
\emph{nearly double} the hallucination rate. BKT-only achieves cheap
mastery gains by rapidly converging $P(L^k)$ on a simulated student
that does not perceive incorrect feedback; a real student exposed to
that feedback acquires misconceptions invisible to BKT-AUC.
AgentRAG-Tutor's CRAG layer trades simulated LG for real-student
safety.

\paragraph{Why the trained PPO picks $d=3$.}
The trained policy converges on $d=3$ across most cold-start mastery
states. This is the analytic optimum of the reward function: the
easy-penalty term $-\lambda\max(0, 3-d)$ vanishes at $d=3$, and the
simulator's difficulty--correctness curve makes $d>3$ produce small
or negative mastery deltas under cold-start.

\paragraph{Multi-domain architecture generalisation.}
The cross-domain reward gap of 1.14 on AI Fundamentals to 1.06 on
Operating Systems is small considering OS has different per-KC priors
and was never seen during PPO training, supporting the claim that the
PPO policy operates on a domain-agnostic abstraction.

\paragraph{Template-disjointness as honest evaluation.}
The intent results show the rule baseline winning under a
template-disjoint split when the same classifier loses under a random
split. Production deployments that report only random-split results
overstate generalisation.

\begin{figure*}[t]
\centering
\includegraphics[width=0.62\textwidth]{BKT.png}
\caption{BKT mastery curves over five sessions for three
representative simulated personas (S04~strong, S08~medium, S19~weak);
five colour-coded curves per persona, one per KC.}
\label{fig:bktcurves}
\end{figure*}

\begin{figure*}[t]
\centering
\includegraphics[width=0.50\textwidth]{PPO.png}
\caption{PPO cumulative reward versus training steps; mean episode
reward rises from about 0.25 at start to 1.14 at 500{,}000 steps;
the smoke-train baseline (0.78 at 50{,}000 steps) is marked as a
dashed point.}
\label{fig:ppocurve}
\end{figure*}

% =====================================================================
\section{Limitations and Ethics}\label{sec:limit}
% =====================================================================

\paragraph{Simulator-bound results.}
The headline tutor numbers in
Tables~\ref{tab:baselines}--\ref{tab:cross} are from a BKT-driven
simulator, not real students. The pilot in Sec.~\ref{sec:pilot}
provides initial external validity; an $N{=}20$ controlled study is
planned for the journal extension.

\paragraph{PPO is inference-only.}
The deployed PPO policy is frozen. Online fine-tuning is left to
future work.

\paragraph{Single-skill BKT.}
The current BKT layer assumes one KC per question. Multi-skill BKT
(q-matrix learning) is not yet wired in.

\paragraph{Difficulty-correctness curve.}
The simulator's $\delta(d)$ mapping is heuristic, not data-fit. A
pilot-data-driven recalibration is part of future work.

\paragraph{Ethics.}
All synthetic student data is clearly labelled as simulated. The
pilot will obtain explicit informed consent, anonymise participants
by SHA-256 of email, and store no identifying information beyond
what is needed for session reconstruction.

% =====================================================================
\section{Conclusion}\label{sec:conclude}
% =====================================================================

We have presented AImentor, a production-grade open-source AI career
mentor that combines multi-provider LLM orchestration, hybrid intent
classification, semantic ATS scoring, and an adaptive tutoring
module --- AgentRAG-Tutor --- that pipelines an agentic CRAG
retrieval loop, a 5-KC paper-calibrated BKT tracker, and a PPO
difficulty agent that consumes the BKT mastery vector as state. On a
20-persona BKT-driven simulator, AgentRAG-Tutor attains an 18.0
percent hallucination rate against 33--34 percent for non-CRAG
baselines while remaining domain-portable: the same trained PPO
checkpoint works on a second curriculum (Operating Systems) with no
fine-tuning. A controlled $N{=}5$ pilot at SRM University AP yields a
$59.7{\pm}20.6$ Hake learning gain and 4.4/5 mean satisfaction,
providing first real-learner external validity for the system.

\begin{thebibliography}{99}

\bibitem{lewis2020rag} P.~Lewis et al., ``Retrieval-augmented
generation for knowledge-intensive NLP tasks,'' in
\emph{Adv. Neural Inf. Process. Syst.} (NeurIPS), 2020.

\bibitem{yan2024crag} S.-Q.~Yan, J.-C.~Gu, Y.~Zhu, and Z.-H.~Ling,
``Corrective retrieval augmented generation,'' arXiv preprint
arXiv:2401.15884, 2024.

\bibitem{asai2024selfrag} A.~Asai, Z.~Wu, Y.~Wang, A.~Sil, and
H.~Hajishirzi, ``Self-RAG: Learning to retrieve, generate, and
critique through self-reflection,'' in \emph{Proc. ICLR}, 2024.

\bibitem{corbett1994bkt} A.~T.~Corbett and J.~R.~Anderson,
``Knowledge tracing: Modeling the acquisition of procedural
knowledge,'' \emph{User Modeling and User-Adapted Interaction},
vol.~4, no.~4, pp.~253--278, 1994.

\bibitem{yudelson2013ibkt} M.~V.~Yudelson, K.~R.~Koedinger, and
G.~J.~Gordon, ``Individualized Bayesian knowledge tracing models,''
in \emph{Proc. AIED}, 2013, pp.~171--180.

\bibitem{piech2015dkt} C.~Piech et al., ``Deep knowledge tracing,''
in \emph{Adv. Neural Inf. Process. Syst.} (NeurIPS), 2015.

\bibitem{schulman2017ppo} J.~Schulman, F.~Wolski, P.~Dhariwal,
A.~Radford, and O.~Klimov, ``Proximal policy optimization
algorithms,'' arXiv preprint arXiv:1707.06347, 2017.

\bibitem{chi2011rl} M.~Chi, K.~VanLehn, D.~Litman, and P.~W.~Jordan,
``An evaluation of pedagogical tutorial tactics for a natural
language tutoring system: A reinforcement learning approach,''
\emph{Int. J. Artif. Intell. Educ.}, vol.~21, no.~1--2,
pp.~83--113, 2011.

\bibitem{linchi2016ibkt} C.~Lin and M.~Chi, ``Intervention-BKT:
Incorporating instructional interventions into Bayesian knowledge
tracing,'' in \emph{Proc. ITS}, 2016, pp.~208--218.

\bibitem{vanlehn2011} K.~VanLehn, ``The relative effectiveness of
human tutoring, intelligent tutoring systems, and other tutoring
systems,'' \emph{Educational Psychologist}, vol.~46, no.~4,
pp.~197--221, 2011.

\bibitem{anderson1995} J.~R.~Anderson, A.~T.~Corbett,
K.~R.~Koedinger, and R.~Pelletier, ``Cognitive tutors: Lessons
learned,'' \emph{J. Learn. Sci.}, vol.~4, no.~2, pp.~167--207,
1995.

\bibitem{khanmigo2023} Khan Academy, ``Khanmigo: Khan Academy's AI
tutor,'' Khan Academy Tech. Rep., 2023.

\bibitem{ragas2023} S.~Es, J.~James, L.~Espinosa-Anke, and
S.~Schockaert, ``RAGAS: Automated evaluation of retrieval augmented
generation,'' arXiv preprint arXiv:2309.15217, 2023.

\bibitem{ouyang2022instructgpt} L.~Ouyang et al., ``Training
language models to follow instructions with human feedback,'' in
\emph{Adv. Neural Inf. Process. Syst.} (NeurIPS), 2022.

\bibitem{pavlik2009pfa} P.~I.~Pavlik~Jr., H.~Cen, and
K.~R.~Koedinger, ``Performance factors analysis: A new alternative
to knowledge tracing,'' in \emph{Proc. AIED}, 2009, pp.~531--538.

\bibitem{karpukhin2020dpr} V.~Karpukhin et al., ``Dense passage
retrieval for open-domain question answering,'' in \emph{Proc.
EMNLP}, 2020.

\bibitem{johnson2019faiss} J.~Johnson, M.~Douze, and H.~Jegou,
``Billion-scale similarity search with GPUs,'' \emph{IEEE Trans.
Big Data}, vol.~7, no.~3, pp.~535--547, 2019.

\bibitem{reimers2019sbert} N.~Reimers and I.~Gurevych,
``Sentence-BERT: Sentence embeddings using Siamese BERT-networks,''
in \emph{Proc. EMNLP}, 2019.

\bibitem{robertson2009bm25} S.~Robertson and H.~Zaragoza, ``The
probabilistic relevance framework: BM25 and beyond,'' \emph{Found.
Trends Inf. Retr.}, vol.~3, no.~4, pp.~333--389, 2009.

\bibitem{onet2023} National Center for O*NET Development, ``O*NET
Online occupational information,'' U.S. Department of Labor, 2023.

\bibitem{srivatsa2023router} K.~C.~Srivatsa et al., ``Harnessing
the power of multiple minds: Lessons learned from LLM routing,''
arXiv preprint, 2023.

\bibitem{ding2024hybrid} D.~Ding et al., ``Hybrid LLM:
Cost-efficient and quality-aware query routing,'' in \emph{Proc.
ICLR}, 2024.

\bibitem{larson2019clinc} S.~Larson et al., ``An evaluation
dataset for intent classification and out-of-scope prediction,''
in \emph{Proc. EMNLP-IJCNLP}, 2019.

\bibitem{casanueva2020banking77} I.~Casanueva et al., ``Efficient
intent detection with dual sentence encoders,'' in \emph{Proc. ACL
Workshop on NLP for Conversational AI}, 2020.

\bibitem{koh2021wilds} P.~W.~Koh et al., ``WILDS: A benchmark of
in-the-wild distribution shifts,'' in \emph{Proc. ICML}, 2021.

\bibitem{singla2021rl4ed} A.~Singla, A.~N.~Rafferty, G.~Radanovic,
and N.~Heffernan, ``Reinforcement learning for education:
Opportunities and challenges,'' in \emph{Proc. EDM}, 2021.

\bibitem{baker2010more} R.~S.~J.~d.~Baker, A.~T.~Corbett,
S.~M.~Gowda et al., ``Contextual slip and prediction of student
performance after use of an intelligent tutor,'' in \emph{User
Modeling, Adaptation, and Personalization}, 2010, pp.~52--63.

\bibitem{rocha2024edurag} E.~Da~Rocha et al., ``Educational
retrieval-augmented generation: A survey,'' arXiv preprint
arXiv:2402.10371, 2024.

\bibitem{khan2023llmtutor} A.~Khan and N.~Heffernan, ``Are large
language models ready to replace intelligent tutoring systems?''
in \emph{Proc. AIED Workshop on Generative AI in Education}, 2023.

\bibitem{rajpurkar2018squad2} P.~Rajpurkar, R.~Jia, and P.~Liang,
``Know what you don't know: Unanswerable questions for SQuAD,''
in \emph{Proc. ACL}, 2018, pp.~784--789.

\bibitem{hake1998} R.~R.~Hake, ``Interactive-engagement versus
traditional methods,'' \emph{Am. J. Phys.}, vol.~66, no.~1,
pp.~64--74, 1998.

\bibitem{kasneci2023chatgpt} E.~Kasneci et al., ``ChatGPT for good?
On opportunities and challenges of large language models for
education,'' \emph{Learning and Individual Differences},
vol.~103, 2023.

\end{thebibliography}

\end{document}